\subsection{公理}

一个逻辑理论是任何这样的理论 \(\mathfrak{G}\) ：其中下面的方案S1至S4提供了隐式公理。

S1. 如果 \(A\) 是 \(\mathfrak{C}\) 中的一个关系，关系 \((A \text{ 或 } A) \Longrightarrow A\) 是 \(\mathfrak{C}\) 的公理\footnote{此方案可以不用字母A或缩写符号 \(\Longrightarrow\) 来表达，如下所述：每当我们有一个关系时，我们通过从左到右书写 \(\vee\), \(\neg\), \(\vee\)，然后给定该关系三次，便得到一个定理。读者可以作为一个练习，以类似方式翻译其他图式的表达式。}.

S2. 若A和B是C中的关系，则关系 \(A \Longrightarrow (A \text{ 或 } B)\) 是C的一个公理。

S3. 若A和B是 \(\mathfrak{G}\) 中的关系，则关系(A或B) \(\Longrightarrow\) (B或A)是G的公理。

S4. 若A、B和C是 \(\mathfrak{C}\) 中的关系，则关系

\[
(A \Rightarrow B) \Rightarrow ((C \text {   或   } A) \Rightarrow (C \text {   或   } B))
\]

是 \(\mathfrak{C}\) 的公理。

这些规则实际上是方案；让我们验证这一点，例如对于S2。设R是通过应用S2得到的一个关系；则存在C中的关系A和B，使得R是关系 \(A \Longrightarrow (A \text{ 或 } B)\) 。设T是C中的一个项，x是一个字母，且 \(A'\) 和 \(B'\) 为关系 \((T|x)A\) 和 \((T|x)B\) ；然后 \((T|x)R\) 与 \(A' \Longrightarrow (A' \text{ 或 } B')\) 相同，因此可以通过应用 S2 得到。

直观上，规则 S1 至 S4 仅仅表达了在通常的数学语言中“或”和“蕴含”这两个词所附带的含义。\footnote{在日常用语中，“或”这个词根据语境有两种不同的含义：当我们用“或”连接两个陈述时，可能意指至少其中一个成立（也可能两者同时成立），或者可能意指两者中仅有一个成立。}.

如果某个逻辑理论 \(\mathfrak{C}\) 是矛盾的，则 \(\mathfrak{C}\) 中的每个关系都是 \(\mathfrak{C}\) 中的定理。因为令 \(A\) 是 \(\mathfrak{C}\) 中的一个关系，使得 \(A\) 与“非 \(A\)”都是 \(\mathfrak{C}\) 中的定理，并且设 \(B\) 是 \(\mathfrak{C}\) 中的任意关系。根据S2，(非 \(A\) ) \(\Longrightarrow\) ((非 \(A\) ) 或 \(B\) ) 是 \(\mathfrak{C}\) 中的定理；因此，根据C1（ \(\S 2\) ，第2条），“(非 \(A\) ) 或 \(B\) ”，也就是说， \(A\Longrightarrow B\) 是 \(\mathfrak{C}\) 中的定理。再次应用C1表明， \(B\) 是 \(\mathfrak{C}\) 中的定理 .

从现在起， \(\mathfrak{C}\) 将表示一个逻辑理论。

